\documentclass[11pt]{article}

% ---------------------------------------------------------------------------
% Kazakhstan Practical Driving Exam Simulator — project report
% Interactive Graphics, Sapienza University of Rome.
% Compile with pdfLaTeX (MiKTeX): pdflatex report.tex  (run twice for the ToC).
% ---------------------------------------------------------------------------

\usepackage[utf8]{inputenc}
\usepackage[T1]{fontenc}
\usepackage{lmodern}   % outline fonts: correct ligatures + copy-pasteable text
\usepackage[a4paper,margin=2.5cm]{geometry}
\usepackage{amsmath,amssymb}
\usepackage{siunitx}
\usepackage{graphicx}
\usepackage{booktabs}
\usepackage{xcolor}
\usepackage{listings}
\usepackage{fancyhdr}
\usepackage{caption}
\usepackage[hidelinks]{hyperref}

% --- Running header: short title left, page number right --------------------
\pagestyle{fancy}
\fancyhf{}
\fancyhead[L]{Driving Exam Simulator}
\fancyhead[R]{\thepage}
\renewcommand{\headrulewidth}{0.4pt}

% --- Code listing style (JavaScript) ---------------------------------------
\definecolor{kw}{rgb}{0.55,0.10,0.45}
\definecolor{cmt}{rgb}{0.40,0.45,0.50}
\definecolor{str}{rgb}{0.10,0.45,0.25}
\definecolor{numcol}{rgb}{0.55,0.55,0.55}
\definecolor{codebg}{rgb}{0.975,0.975,0.975}

\lstdefinelanguage{JS}{
  keywords={const,let,var,function,return,if,else,for,of,in,new,this,class,
            export,import,from,true,false,null,undefined,get,set,Math},
  morecomment=[l]{//},
  morecomment=[s]{/*}{*/},
  morestring=[b]',
  morestring=[b]",
  sensitive=true,
}
\lstset{
  language=JS,
  basicstyle=\ttfamily\footnotesize,
  keywordstyle=\bfseries\color{kw},
  commentstyle=\itshape\color{cmt},
  stringstyle=\color{str},
  numbers=left,
  numberstyle=\tiny\color{numcol},
  numbersep=9pt,
  backgroundcolor=\color{codebg},
  frame=single,
  rulecolor=\color{black!25},
  breaklines=true,
  showstringspaces=false,
  captionpos=b,
  tabsize=2,
  aboveskip=1.1em,
  belowskip=0.6em,
}
\renewcommand{\lstlistingname}{Listing}

\newcommand{\code}[1]{\texttt{#1}}

% ===========================================================================
\begin{document}

% --- Title page ------------------------------------------------------------
\begin{titlepage}
  \centering
  \vspace*{3cm}
  {\Huge\bfseries Kazakhstan Practical\\[2mm] Driving Exam Simulator\par}
  \vspace{4mm}
  {\large Category~B - Astana Autodrome\par}
  \vspace{2cm}

  % TODO: replace with the real team member names and student emails.
  \begin{tabular}{cc}
    \large Damir Tassybayev \\
    \texttt{tassybayev.kostanay@gmail.com} \\[4mm]
  \end{tabular}

  \vspace{2cm}
  {Course: Interactive Graphics\par}
  {Instructor: Prof.\ Marco Schaerf\par}
  {Sapienza University of Rome\par}

  \vfill
  % To add a cover image: \includegraphics[width=0.6\textwidth]{figures/cover.png}
\end{titlepage}

\thispagestyle{fancy}
\tableofcontents
\newpage

% ===========================================================================
\section{Introduction}

This project is a browser-based 3D simulation of the \textbf{practical driving
examination of the Republic of Kazakhstan} for the category-B licence. The
player drives a car around a faithful recreation of an Astana \emph{autodrome}
(driving-test ground) and is graded, in real time, by an automated examiner that
reproduces the official penalty system: turn-signal checkpoints, a hill start
(\emph{estakada}), traffic lights, boundary lines, and two timed parking
exercises. The simulator was developed as the final project for the Interactive
Graphics course at Sapienza University of Rome.

\subsection{Graphics environment}

The application runs in the browser through WebGL, using \textbf{Three.js}
(release r160) as the rendering layer and \textbf{tween.js} for eased
interpolation. The renderer is a \code{WebGLRenderer} with antialiasing enabled,
the \code{sRGB} output colour space, soft shadows (\code{PCFSoftShadowMap}), and
a device pixel ratio capped at~2 to avoid over-drawing on high-DPI screens.

A deliberate design choice is that the project has \textbf{no build step}. All
modules are native ES modules resolved through an \code{importmap} in
\code{index.html}, and every third-party library is vendored locally under
\code{lib/}. This keeps the repository self-contained and lets it run directly
from GitHub Pages (or any static file server) without bundling.

\subsection{Running the simulator}

Because the project uses ES modules, it must be served over HTTP rather than
opened from the file system. Any static server works, for example:

\begin{lstlisting}[language={},caption={Running the simulator locally.}]
python -m http.server 8000     # from the repository root
# then open http://localhost:8000
\end{lstlisting}

It has been tested on Chromium-based browsers and Firefox on Windows. Because the
controls are keyboard-driven, mobile devices are not supported.

% ===========================================================================
\section{User Manual}

\subsection{Premise and objective}

The candidate starts the car at the \textsc{start} box at the top of the
autodrome and drives a fixed route through a sequence of exercises. The automated
examiner accumulates \emph{penalty points} for mistakes. The examination is
\textbf{passed} if, and only if, the candidate completes the course and the
cumulative penalty stays \emph{below 100 points}; reaching 100 points, whether
through one critical error or several smaller ones, ends the examination
immediately as a failure. Penalties fall into three severity tiers, summarised in
Table~\ref{tab:tiers}.

\begin{table}[h]
  \centering
  \caption{Penalty severity tiers.}
  \label{tab:tiers}
  \begin{tabular}{rll}
    \toprule
    Points & Severity & Examples \\
    \midrule
    100 & Critical (instant fail) & Crossing on a red/yellow light; not finishing a parking in time \\
    25  & Major violation & Missing a full stop; heavy hill rollback \\
    5   & Minor mistake & Missing a turn signal at a checkpoint \\
    20  & Boundary touch & A wheel on a penalty (boundary) line \\
    \bottomrule
  \end{tabular}
\end{table}

\subsection{Controls}

The car is driven with the keyboard; the mouse orbits the camera only in the
free-orbit view. The full scheme is listed in Table~\ref{tab:controls}.

\begin{table}[h]
  \centering
  \caption{Controls.}
  \label{tab:controls}
  \begin{tabular}{ll}
    \toprule
    Input & Action \\
    \midrule
    \texttt{W} / \texttt{S} or $\uparrow$ / $\downarrow$ & Accelerate / brake and reverse \\
    \texttt{A} / \texttt{D} or $\leftarrow$ / $\rightarrow$ & Steer left / right \\
    \texttt{Space} & Full-stop handbrake (toggle) \\
    \texttt{Q} / \texttt{E} & Left / right turn signal (toggle) \\
    \texttt{H} & Hazard lights (toggle) \\
    \texttt{C} & Cycle the camera view \\
    \texttt{L} & Toggle the headlights \\
    \texttt{O} & Open / close the doors \\
    \texttt{R} & Reset the car and the examination \\
    Mouse drag / wheel & Orbit / zoom (in the \emph{orbit} camera only) \\
    \bottomrule
  \end{tabular}
\end{table}

\subsection{Camera modes}

The \texttt{C} key cycles four views, handled by the \code{CameraManager}. The
default \emph{fixed} view is rigidly bolted to the car: it translates and rotates
with the body so the car always sits in the same place on screen. \emph{Cockpit}
places the camera at the driver's seat looking forward. \emph{Top} looks straight
down, useful for the parking exercises. \emph{Orbit} frees the camera with
\code{OrbitControls} so the player can drag and zoom around the car. The
non-orbit modes are positioned each frame from a car-local offset converted to
world space, so they follow the body exactly.

\subsection{HUD overview}

Two panels overlay the scene. The top-left panel reports the live vehicle state:
speed in km/h, gear (D, R or P when the handbrake is engaged), the active camera,
the headlight state, and the current turn signal. The top-right \emph{examiner}
panel shows the running penalty total, colour-coded from green to amber to red as
it climbs, together with the most recent violation. When the examination ends, a
full-screen overlay announces \textsc{passed} or \textsc{failed} with the final
penalty total and, on failure, the reason.

% ===========================================================================
\section{The Car: hierarchical model}

The car is the project's complex hierarchical model and is built entirely in code
in \code{world/car.js}; nothing about it is imported, in line with the course rule
that animations may not be imported.

\subsection{Part hierarchy and assembly}

The car is a tree of \code{THREE.Group} nodes. The root node carries the whole
car and is the only node the driving model moves; every child therefore follows
automatically. The hierarchy is:

\begin{itemize}
  \item \textbf{root} \textemdash{} position and heading set by the driving model;
  \begin{itemize}
    \item \textbf{body} \textemdash{} lower body, cabin, glass band, chrome
      bumpers, seats; the emissive head-, brake- and indicator-light lenses; the
      two hinged \emph{doors}; and the \emph{steering wheel} (a torus nested in
      two helper groups so it can spin about its own axle regardless of its tilt);
    \item \textbf{front-left} and \textbf{front-right pivots} \textemdash{} each a
      group that yaws for steering and contains a wheel that spins for rolling;
    \item \textbf{rear-left} and \textbf{rear-right wheels} \textemdash{} spin only;
    \item \textbf{two head\-light spot lights} aimed forward.
  \end{itemize}
\end{itemize}

The car's local forward direction is $+X$ and up is $+Y$. Key dimensions are a
half-wheelbase of \SI{1.45}{m}, a half-track of \SI{0.82}{m}, a wheel radius of
\SI{0.42}{m} and a steering lock of \code{MAX\_STEER} $= 0.52$ rad ($\approx
30^\circ$). Because the front wheels are children of their steering pivots, they
turn in place and still roll correctly; this two-level joint/visual split is what
lets every joint rotate independently without disturbing the others.

\subsection{Materials and lights}

The body uses \code{MeshStandardMaterial} instances (metallic paint, dark glass,
chrome bumpers, rubber tyres). The light lenses use emissive materials whose
intensity is animated: the headlight lenses and two non-shadowing
\code{SpotLight}s fade in when the lights are toggled, the brake lenses brighten
while braking or reversing, and the amber indicator lenses blink at about
\SI{2}{Hz}. Car-local \emph{left} is the $-Z$ side and \emph{right} is $+Z$, so
the left/right indicators light the correct side of the body.

% ===========================================================================
\section{The Autodrome (environment)}

\subsection{Map image and pixel-to-world mapping}

The driving ground is a single photographic image of a real Astana autodrome,
laid flat on a textured plane (\code{world/examGround.js}). The image
($1237\times848$~px) is loaded asynchronously with \code{TextureLoader} and the
ground plane is sized to the \emph{same aspect ratio}, so a pixel on the image
maps exactly to a point on the ground. The plane is oriented so that the image
top is world north ($-Z$) and the image right is world east ($+X$).

Every element of the examination is positioned by reading its location directly
off the image in pixels and converting it with the helper in
\code{world/mapCoords.js}:

\begin{lstlisting}[caption={Pixel-to-world mapping (\code{mapCoords.js}).}]
export function px2world(px, py) {
  return {
    x: (px / IMG_W - 0.5) * AUTO_W,   // 0..IMG_W  (left->right)  -> world X
    z: (py / IMG_H - 0.5) * AUTO_H,   // 0..IMG_H  (top->bottom)  -> world Z
  };
}
\end{lstlisting}

This single source of truth keeps "what is painted on the image" and "what the
car interacts with" in agreement: the car start, the hill, and all examination
lines are expressed in image pixels and converted through \code{px2world}.

\subsection{The overpass (estakada)}

The hill-start exercise is a real raised structure built in \code{world/overpass.js}.
Its footprint is taken from image pixels (it spans image $x=550$ on the east,
where the west-bound car arrives, to $x=210$ on the west). The deck is an
extruded side profile: an up-ramp, a flat crest carrying a painted \textsc{stop}
line and a \textsc{stop} sign on the right edge, and a down-ramp. A height
function returns the deck elevation at any point, and zero elsewhere:

\begin{lstlisting}[caption={Overpass deck height (\code{overpass.js}).}]
function heightAt(x, z) {
  if (z < z0 || z > z1 || x > xA || x < xD) return 0;
  if (x >= xB) return H * (xA - x) / (xA - xB); // up ramp
  if (x <= xC) return H * (x - xD) / (xC - xD); // down ramp
  return H;                                     // crest
}
\end{lstlisting}

The driving model samples this function so the car physically rides up, over and
down the deck (Section~\ref{sec:terrain}).

\subsection{Course markings and traffic lights}

The examination lines that must be visible are drawn on the ground by
\code{world/courseMarks.js}: the boundary ``penalty'' lines as thin black strips,
and a red stop line at each traffic-light location. At each traffic light a small
hierarchical light model (post, housing and three lamps) is placed and driven by a
green\,$\to$\,yellow\,$\to$\,red cycle, desynchronised between lights. The
remaining rule lines (turn-signal, full-stop, parking, finish) are invisible
triggers by default to keep the photographic map clean; a debug flag draws them
all for verification.

\subsection{Lighting and shadows}

The scene is lit by a \code{HemisphereLight} ambient term, so no surface is fully
black, and a \code{DirectionalLight} acting as the sun, which casts the scene
shadows through a $2048\times2048$ shadow map. The car's headlights add two
forward spot lights at night. Shadows use percentage-closer filtering for soft
edges, and the renderer works in the sRGB colour space for correct tone.

% ===========================================================================
\section{Driving model}
\label{sec:terrain}

\subsection{Kinematic bicycle model}

The car is driven by a kinematic \emph{bicycle model} (\code{vehicle.js}); no
physics engine is used, keeping the motion light and predictable. The state is
the planar position $(x,z)$, the heading $\psi$, the signed speed $v$, and the
current steering angle $\delta$. Throttle accelerates, braking decelerates and
then engages reverse, and releasing both applies rolling friction toward zero.
With wheelbase $L=\SI{2.9}{m}$ the pose integrates as

\[
  \dot{\psi} = \frac{v}{L}\tan\delta,
  \qquad
  \dot{x} = v\cos\psi,
  \qquad
  \dot{z} = -\,v\sin\psi .
\]

Speed is clamped to $[-5,\,15]\ \mathrm{m/s}$ ($\approx -18$ to $54$~km/h) and the
steering lock is reduced at higher speed for stability. The travelled distance is
fed back to the wheels so the rolling animation always matches the real speed.

\begin{lstlisting}[caption={Pose integration (\code{vehicle.js}).}]
if (Math.abs(this.speed) > 1e-3) {
  this.heading += (this.speed / WHEELBASE) * Math.tan(this.steer) * dt;
}
this._dir.set(Math.cos(this.heading), 0, -Math.sin(this.heading));
const dist = this.speed * dt;
this.x += this._dir.x * dist;
this.z += this._dir.z * dist;
\end{lstlisting}

\subsection{Riding the terrain}

When the car is over the overpass, the model raises it to the deck height and
pitches it to the local slope. The pitch is taken from two height samples a short
distance $e$ ahead of and behind the car along its travel direction, and a
component of gravity along the slope is added to the speed, so the hill must be
actively held:

\[
  \theta_{\mathrm{pitch}} = \operatorname{atan2}\!\big(h_f - h_b,\ 2e\big),
  \qquad
  \dot v \mathrel{+}= -\,g\,\sin\theta_{\mathrm{pitch}} .
\]

The resulting position and orientation are pushed to the car root as a quaternion
(yaw then pitch), so the body sits correctly on the ramp.

\subsection{Handbrake and stall}

Pressing \texttt{Space} toggles a full-stop handbrake: the speed is forced to zero
and throttle is ignored until it is released, with the brake lights held on. The
model also detects a stall when throttle and brake are held together near a
standstill, which is reported to the examiner.

% ===========================================================================
\section{Animations}

All animation is computed in JavaScript every frame; nothing is imported.

\begin{itemize}
  \item \textbf{Wheels.} All four wheels spin about their axle by an angle that
    advances with the travelled distance, $\Delta\theta = \dfrac{\mathrm{d}s}{r}$,
    where $r$ is the wheel radius. The two front wheels additionally yaw with the
    steering pivots.
  \item \textbf{Steering wheel.} The cabin steering wheel rotates about its own
    axle proportionally to the steering angle.
  \item \textbf{Doors.} Each door is a hinge group; opening and closing is driven
    by \textbf{tween.js} with a \code{Quadratic.Out} easing for a smooth swing.
  \item \textbf{Lights.} Headlight and brake intensities ease toward their target;
    the turn indicators blink; the brake lenses respond to braking and reversing.
  \item \textbf{Traffic lights.} Each light advances through its
    green/yellow/red cycle and updates the emissive lamp intensities.
\end{itemize}

% ===========================================================================
\section{The examiner: scoring system}

The examination logic is a generic, data-driven engine that reads a course
definition and applies a rule per element each frame, accumulating penalties.

\subsection{Scoring engine}

The \code{Scoring} class (\code{exam/scoring.js}) keeps a running total and an
ordered event log, and fails the examination as soon as a critical (100-point)
penalty fires or the total reaches the threshold:

\begin{lstlisting}[caption={Penalty accumulation (\code{scoring.js}).}]
penalize(exercise, rule, points, t = 0) {
  if (this.isOver) return;
  this.total += points;
  this.log.push({ t, exercise, rule, points });
  this.lastViolation = this.log[this.log.length - 1];
  if (points >= FAIL_THRESHOLD || this.total >= FAIL_THRESHOLD) {
    this.status = 'failed';
    this.failReason = rule;
  }
}
\end{lstlisting}

\subsection{Data-driven course}

The whole route lives in one editable configuration file,
\code{exam/course.js}, with every position expressed in image pixels. It groups
the elements by type: \emph{signal lines} (a crossing requires a given indicator),
\emph{light lines} (a crossing on red/yellow fails), a \emph{stop line} (a full
stop is required before crossing), \emph{penalty lines} (a wheel on the line
costs points), two timed \emph{parkings}, and the \emph{finish line}. The engine
(\code{exam/exam.js}) converts each pixel segment to a world segment once at
start-up and then evaluates it generically.

\subsection{Line-crossing detection}

The car is reduced to a reference point (its front axle) and its four wheels in
world space. A line is the segment $A\!\to\!B$; the side of the car relative to
the line is the sign of the 2-D cross product

\[
  s = (B_x - A_x)(P_z - A_z) - (B_z - A_z)(P_x - A_x).
\]

A crossing is registered when $s$ flips sign between frames \emph{and} the
projection of $P$ onto the segment lies within its extent. Each rule then reacts:
a signal line checks the active indicator, a light line reads its traffic
signal, the stop line checks that a one-second stop was registered nearby, and the
finish line ends the examination. The penalty (boundary) lines instead use a
point-to-segment distance against each wheel, with a short cooldown so a wheel
resting on a line is charged once rather than every frame.

\subsection{Parking and hill exercises}

Each parking exercise arms when the car crosses its trigger line and then gives a
fixed time budget (30~s). It succeeds when the required wheels (the two rear
wheels for the $90^\circ$ box, the two right wheels for parallel parking) are held
stationary inside the target line for one second; if the budget expires first, the
examination fails. The hill exercise is scored against the overpass footprint:
rolling back more than \SI{30}{cm} from the furthest-forward point, or stopping
more than \SI{0.5}{m} from the \textsc{stop} line (or not stopping at all), each
costs a major penalty.

% ===========================================================================
\section{Resource attribution}

\paragraph{Libraries (not developed by the team).}
\begin{itemize}
  \item \textbf{Three.js} (MIT) \textemdash{} \url{https://threejs.org/}; release
    r160, used for rendering, the scene graph and the \code{OrbitControls} addon.
  \item \textbf{tween.js} (MIT) \textemdash{}
    \url{https://github.com/tweenjs/tween.js/}; eased interpolation for the door
    animation.
\end{itemize}
Both libraries are vendored under \code{lib/} so the repository is self-contained.

\paragraph{Assets.}
The autodrome ground is a single photographic top-down image of an Astana
driving-test ground (\code{assets/textures/autodrome.png}). All other geometry
(the car, the overpass, the traffic lights, the course markings) and all
animation are generated in code by the team; no 3D models or animation clips are
imported.

\paragraph{Code.}
All application code under \code{src/} was written by the team. The procedural
asphalt textures are generated at run time on an HTML canvas.

\end{document}
